\subsection{Metodología}
Definido el problema y su estado de conflicto, corresponde establecer una metodología que permita abordarlo de manera estructurada y que al mismo tiempo se adapte a las características de un problema complejo con múltiples actores, restricciones normativas y condiciones de uso exigentes.\\
\subsubsection*{Metodología Doble Diamante} ~
Se propone utilizar la metodología del Doble Diamante \cite{ref_design_council_double_diamond} porque permite organizar el proceso de diseño en etapas divergentes y convergentes, favoreciendo primero una comprensión amplia del problema y luego una definición más precisa de los criterios de intervención. Esta metodología evita pasar prematuramente a la solución y prioriza la exploración del contexto, del usuario y del conflicto antes de formular respuestas. Es especialmente adecuada para este proyecto porque el problema del traspaso clínico no se resuelve simplemente con una aplicación, sino que requiere comprender primero qué necesita realmente el profesional asistencial y cuáles son las restricciones reales del entorno hospitalario.\\
\begin{itemize}
    \item \textbf{Etapa Discover:} En esta etapa se busca comprender el contexto hospitalario interno, identificar los actores del flujo asistencial, reconocer los puntos de quiebre en el traspaso de información entre turnos y levantar las necesidades, tensiones y restricciones reales del sistema. Aquí se sitúan herramientas como el mapa de actores, el buyer persona, el mapa de empatía y el modelado del flujo AS-IS.
    \item \textbf{Etapa Define:} En esta fase se sintetiza la información levantada para delimitar con precisión el problema de diseño. Se define el estado de conflicto, las necesidades críticas del usuario principal, los criterios de no interferencia con la labor asistencial y los requerimientos que deberá cumplir la futura solución para ser adoptada en un entorno de alta carga operativa.
    \item \textbf{Etapa Develop:} En esta etapa se proyectan líneas de solución preliminares, considerando distintas formas de capturar el contexto clínico cotidiano, organizarlo cronológicamente y hacerlo disponible para el equipo entrante. Se formulan conceptos parciales y criterios de interacción acordes al contexto de uso real del personal asistencial.
    \item \textbf{Etapa Deliver:} En el diseño preliminar documentado, esta etapa se plantea como la consolidación de lineamientos conceptuales y requisitos de diseño que orienten el desarrollo posterior del sistema y de sus interfaces; no constituye evidencia de implementación ni de un prototipo final.
\end{itemize}
